% Source PDF: source/普通高考/2014/2014广东理.pdf
% Page: 1; PDF embedded bitmap attempt: no image objects were present (vector source).
% Grid evidence: tmp/repaint_2014_guangdong_science/page1_grid100.jpg and q6_block_grid50.jpg.
% Original-side crop: tmp/repaint_2014_guangdong_science/q6_bar_original.png, from the no-grid page render.
% Elements: near-sighted-rate y-axis (percent), grade x-axis with arrow, three unit-width bars,
% dashed horizontal guides at 10%, 30%, and 50%, and the origin/figure labels.
% Relations: the bars for 小学, 初中, 高中 have heights 10%, 30%, 50%; each dashed guide
% ends at the corresponding bar's right boundary. The x-axis categories are ordered by school stage.
% solid segments: bar outlines, x/y axes, axis arrowheads, and category tick marks.
% dashed segments: horizontal reference guides from the y-axis to the matching bar boundary.
% labels: 近视率/% above the y-axis, 年级 to the right of the x-axis, O at the origin, 图2 below.
\begin{tikzpicture}[baseline]
\begin{axis}[
    width=6.0cm,
    height=3.6cm,
    axis lines=none,
    clip=false,
    xmin=-0.65,
    xmax=4.55,
    ymin=-21,
    ymax=57,
    ticks=none,
]
    % Dashed percentage guides terminate at the matching bar boundary.
    \draw[dashed,line width=0.55pt] (axis cs:0,10) -- (axis cs:1,10);
    \draw[dashed,line width=0.55pt] (axis cs:0,30) -- (axis cs:2,30);
    \draw[dashed,line width=0.55pt] (axis cs:0,50) -- (axis cs:3,50);

    % Three grade bars: primary, junior high, senior high.
    \addplot[
        ybar,
        bar width=0.62,
        draw=black,
        fill=white,
        line width=0.7pt,
    ] coordinates {
        (1,10)
        (2,30)
        (3,50)
    };

    % Axes and arrowheads.
    \draw[->,line width=0.8pt] (axis cs:0,0) -- (axis cs:4.05,0);
    \draw[->,line width=0.8pt] (axis cs:0,0) -- (axis cs:0,55);

    % Axis labels and origin.
    \node[anchor=south west,inner sep=0pt,font=\normalsize] at (axis cs:0.12,53.5) {近视率/\%};
    \node[anchor=west,inner sep=0pt,font=\normalsize] at (axis cs:4.15,1.2) {年级};
    \node[anchor=north east,inner sep=0pt,font=\normalsize] at (axis cs:-0.08,-0.5) {O};

    % Category ticks and labels.
    \draw[line width=0.45pt] (axis cs:1,0) -- (axis cs:1,-1.2);
    \draw[line width=0.45pt] (axis cs:2,0) -- (axis cs:2,-1.2);
    \draw[line width=0.45pt] (axis cs:3,0) -- (axis cs:3,-1.2);
    \node[anchor=north,inner sep=0pt,font=\normalsize,rotate=-45] at (axis cs:1,-2.0) {小学};
    \node[anchor=north,inner sep=0pt,font=\normalsize,rotate=-45] at (axis cs:2,-2.0) {初中};
    \node[anchor=north,inner sep=0pt,font=\normalsize,rotate=-45] at (axis cs:3,-2.0) {高中};

    % Percentage ticks and labels.
    \draw[line width=0.45pt] (axis cs:0,10) -- (axis cs:-0.08,10);
    \draw[line width=0.45pt] (axis cs:0,30) -- (axis cs:-0.08,30);
    \draw[line width=0.45pt] (axis cs:0,50) -- (axis cs:-0.08,50);
    \node[anchor=east,inner sep=0pt,font=\normalsize] at (axis cs:-0.12,10) {10};
    \node[anchor=east,inner sep=0pt,font=\normalsize] at (axis cs:-0.12,30) {30};
    \node[anchor=east,inner sep=0pt,font=\normalsize] at (axis cs:-0.12,50) {50};

    \node[anchor=north,inner sep=0pt,font=\normalsize] at (axis cs:2,-16.5) {图2};
\end{axis}
\end{tikzpicture}
